\documentclass[12pt]{article}
\usepackage{amsmath,amssymb,latexsym}
%\input picmacs
%\input BoxedEPS
%\HideDisplacementBoxes
%\SetTexturesEPSFSpecial

\topmargin=-.25in
 \textheight=8.2in

\begin{document}

\begin{center} {\bf  \Large{ Inequalities for the Perimeter of an Ellipse}} 
\end{center}

\vspace{24pt}



\begin{center}{Roger W. Barnard, Kent Pearce, Lawrence Schovanec\\ 
Department of Mathematics and Statistics\\ Texas Tech University }
\end{center}


%\footnote{ Supported
%by NSF grant ECS-9720357  and  Texas Advanced Research Program Grant No.\
%003644-123.}

\vspace{18pt} 
\noindent{\bf{1. Introduction}}
\vspace{6pt}

Let $a,b$ denote the lengths of the  semi-major and -minor axes of an ellipse 
with
eccentricity $e=(1/a)\sqrt{a^2 - b^2}$  and whose perimeter is $L(a,b)$ . 
The
problem of approximating  $L(a,b)$ is an ancient one.  An excellent 
account of this
problem is found in a  {\sc monthly} article   by
G. Almkvist and B.  Berndt \cite{AB}. In the article \cite{AB},
several  approximations that have been proposed over a period of nearly 4 
centuries
are presented. The  approximation for $L(a,b)/(\pi (a+b))$
$$
A(a,b)=\frac{2}{a+b}\left(\frac{a^{3/2}+b^{3/2}}{2}\right)^{2/3}
$$ 
was given by Muir
in 1883.    In 1996,   Matti Vuorinen in the paper ``Hypergeometric 
Functions in
Geometric Function Theory'' \cite{MV}, raised the question, 

$$
\begin{array}{llr}
\hspace{0.1in}& \mbox{Is it true that } A(a,b) \mbox{ is an approximation to }
L(a,b)/(\pi (a+b)) & \mbox{(Q1)} \\
&\mbox{from {\em below} throughout the entire range of eccentricity } e \mbox{?}   
\end{array}
$$

As is often the case for a mathematical conjecture, the insights  gained 
in the course of its
resolution  and   extended ramifications of the result supersede the 
original question. In this paper
we develop a technique to answer  question  (Q1) that readily extends 
to a broader class of
inequalities, all of which are motivated by the problem of approximating 
elliptical perimeter.  The
verification of each inequality is accomplished  by showing the 
positivity of an infinite  series.  
The proof of the positivity of the  infinite series is achieved by 
utilizing a computer algebra system
to execute a Sturm sequence argument.  
 

In verifying (Q1) it will be shown that $A(a,b)$ exhibits a monotonic 
dependence on eccentricity. As a
consequence of this  property,  bounds on the error, $E=L(a,b)/(\pi (a+b))-A(a,b)$, 
are readily
established.  As pointed out in \cite{AB}, a primary motivation for  
calculating  elliptical
perimeter arose from astronomical considerations. This
article will illustrate how
 bounds on the error $E$ may be utilized to determine the accuracy of  
 approximations to elliptical perimeter in calculating planetary orbits.

\vspace{18pt} 
\noindent{\bf{2. Approximations of Elliptic Perimeter}}
\vspace{12pt}
  
Recall that  if an ellipse is described by the parametric equations
$x=a\cos\phi$ and $y=b\sin \phi, \; 0 \leq \phi \leq 2 \pi$, then  the 
perimeter $L$
of the ellipse  is given by
$$ 
 L=L(a,b)=\int_0^{2\pi} \sqrt{(a^2 \sin^2 \phi + b^2 \cos^2 \phi)} \; 
d\phi. 
$$  
The perimeter of the ellipse can be expressed exactly in terms of Gauss's 
ordinary hypergeometric series
$$ 
F(a,b; c: x) =\sum_{k=0}^{\infty} \frac {(a)_k (b)_k}{(c)_k k!} x^k, \;\; 
|x|<1
$$ 
where $a,b,$ and $c$ denote arbitrary complex numbers and $(\alpha)_k$ is
defined  by
$$
 (\alpha)_k= \alpha (\alpha +1) (\alpha +2) \cdots (\alpha +k-1).
$$
The following  expressions for the elliptical perimeter are due, respectively,  to
Maclaurin (1742), Euler (1773),  and Ivory (1796) (see \cite {AB} and 
references
therein) and  illustrate the  connection between elliptical perimeter and the 
hypergeometric functions.
 
\vspace{12pt}
\noindent {\em Proposition.  Let $x=a\cos\phi$ and $y=b\sin \phi, \;\; 
0\leq
\phi \leq 2 \pi$ and let $e=(1/a)\sqrt{a^2 - b^2}$ be the eccentricity of 
the ellipse. 
Then
$$
\begin{array}{rcl} L(a,b) &=&2\pi a F\left(\frac{1}{2}, -\frac{1}{2}; 1: 
e^2\right)\\
[2ex] & = & \pi \sqrt{2(a^2 + b^2) }\; F\left(-\frac{1}{4}, \frac{1}{4}; 
1: 
\left (\frac{  a^2 -b^2}{a^2 + b^2} \right)^2\right) \\ [2ex] 
 &=& \pi (a+b) F(-\frac{1}{2},- \frac{1}{2}; 1: \lambda^2)
\end{array}
$$ 
where}
$$
\lambda=\frac{a-b}{a+b}.
$$ 


The approximations presented in \cite{AB} are reproduced in the list 
below.  In the second column
of the table the approximation $ A(\lambda)$ to  the  `normalized' length 
$ L(a,b)/[\pi (a+b)]$ is
given.  The third column provides the first nonzero term in the power 
series for the difference 
$L(a,b)/[\pi (a+b)] - A(\lambda)$.

\noindent Table 1.
\footnotesize
$$
\begin{array}{llr}
 & \hspace{.5in} A(\lambda) &  \displaystyle{\frac{L(a,b)}{ \pi
(a+b)} - A(\lambda)}
\\    \hline  \hline \\
\mbox{Kepler (1609)} \hspace{.5in}  &
 \displaystyle{\frac{2\sqrt{ab}}{a+b}=\sqrt{1-\lambda^2}} &  
\displaystyle{\frac{3}{4}\lambda^2} 
\\ [4ex] 
\mbox{Euler (1773)} \hspace{.5in}  &
\displaystyle{\frac{\sqrt{2(a^2+b^2)}}{a+b}= \sqrt{1+\lambda^2} } & 
\displaystyle{-\frac{1}{4}\lambda^2} 
\\ [4ex] 
\mbox{Sipos (1792)} \hspace{.5in}   &
\displaystyle{\frac{2(a+b)}{(\sqrt{a}+\sqrt{b})^2}= \frac{2}{1 + 
\sqrt{1-\lambda^2}
}} & 
\displaystyle{-\frac{7}{64}\lambda^4} 
\\ [4ex] 
\mbox{Peano (1889)} \hspace{.5in}   &
\displaystyle{ \frac{3}{2}-\frac{\sqrt{ab}}{a+b}=
\frac{3}{2}-\frac{1}{2}\sqrt{1-\lambda^2} } & 
\displaystyle{-\frac{3}{64}\lambda^4}\\ [4ex]
 \mbox{Muir (1883)}  &
\displaystyle{\frac{2}{a+b}\left(\frac{a^{3/2}+b^{3/2}}{2}\right)^{2/3}} 
& 
\displaystyle{\frac{1}{64}\lambda^4}
\\ [3ex] &
\displaystyle{=\frac{1}{2^{2/3}}((1+\lambda)^{\frac{3}{2}}+(1-\lambda)^{\frac{3}{2}}})^{\frac{2}{3}}
&
\\ [4ex] 
\mbox{Lindner (1904)}  & 
\displaystyle{\left\{ 1+\frac{1}{8} 
\left(\frac{a-b}{a+b}\right)^2\right\}^2} &
\displaystyle{ \frac{1}{2^8}\lambda^6}
\\ [3ex]
 & \displaystyle{=\left(1+\frac{1}{8}\lambda^2\right)^2}  & \\
 & & \\  \mbox{Ramanujan  (1914 a)} &
\displaystyle{\frac{3-\sqrt{(a+3b)(3a+b)}}{a+b}} & 
\displaystyle{ \frac{1}{2^9}\lambda^6}
\\  [2ex] &
\displaystyle{=3-\sqrt{4-\lambda^2}}  &
\\ [2ex] 
 \mbox{Ramanujan (1914 b)}  &  
\displaystyle{1 +\frac{3 \left((a-b)(a+b)\right) ^2}{10 +
\sqrt{4-3\left((a-b)/(a+b)\right) ^2}}} & 
\displaystyle{\frac{3}{2^{17}}\lambda^{10}}
\\ [4ex]
 & \displaystyle{=1+\frac{3\lambda^2}{10+\sqrt{4-3\lambda^2}}}  &
\\ [4ex] 
& \displaystyle{=2 \frac{ ( 1 +  \sqrt{1-\lambda^2})^2 + \lambda^2
\sqrt{1-\lambda^2}}  {( 1 +  \sqrt{1-\lambda^2}) 
(1+\sqrt[4]{1-\lambda^2})^2}}  &
\\ [4ex]
 \mbox{Bronshtein (1964)} &
\displaystyle{\frac{1}{16}\frac{64(a+b)^4-3(a-b)^4}{(a+b)^2(3a+b)(a+3b)}} 
& 
\displaystyle{\frac{9}{2^{14}}\lambda^{8}}
\\ [3ex] 
& \displaystyle{=\frac{ 64-3\lambda^4}  {64-16\lambda^2}}  &
\\ [2ex] \hline
\end{array}
$$

% *** Table continued
\noindent Table 1. (continued)
$$
\begin{array}{llr}
 & \hspace{.5in} A(\lambda) &  \displaystyle{\frac{L(a,b)}{ \pi
(a+b)} - A(\lambda)}\\
  \hline  \hline \\
 \mbox{Selmer  (1975,a)} &
\displaystyle{1+\frac{4(a-b)^2}{(5a+3b)(3a+5b)}} & 
\displaystyle{\frac{3}{2^{10}}\lambda^{6}}
\\ [3ex] &
\displaystyle{=1 + \frac{\lambda^2}{4}\frac{16}{16-\lambda^2}}  &
\\ [4ex]  
 \mbox{Selmer  (1975,b)} & 
\displaystyle{\frac{1}{8}\left(12+\left(\frac{a-b}{a+b}
\right)^2-\frac{2\sqrt{2(a^2+6ab+b^2)}}{a+b}\right)}  & 
\displaystyle{\frac{5}{2^{14}} \lambda^{8}}
\\  [3ex]
& \displaystyle =\frac{3}{2} + \frac{1}{8}
\lambda^2-\frac{1}{2}\sqrt{1-\frac{1}{2}\lambda^2} &
\\  [4ex]
\mbox{Almkvist (1978)} &
\displaystyle{2\frac{2(a+b)^2-(\sqrt{a}-\sqrt{b})^4}{(a+b)[(\sqrt{a}+\sqrt{
b})^2 +
2\sqrt{2}\sqrt{a+b}\sqrt[4]{ab}]}} & 
\displaystyle{ -\frac{15}{2^{14}}\lambda^{8}}
\\ [4ex]  
\mbox{Jacobsen (1985)} &
\displaystyle{ 
\frac{256-48\lambda^2-21\lambda^4}{256-112\lambda^2+3\lambda^4}} & 
\displaystyle{\frac{33}{2^{18}}\lambda^{10}} 
\\ [2ex] \hline
\end{array}
$$

\normalsize

\vspace{12pt}
Let us now reconsider question (Q1). If we utilize Maclaurin's exact 
expression for elliptical perimeter given
in  the {\em Proposition},  then  Muir's approximation  amounts to the 
statement that 
$$  
2\pi a F\left(\frac{1}{2}, -\frac{1}{2}; 1: e^2\right) \approx 2\pi
\left(\frac{a^{3/2}+b^{3/2}}{2} \right)^{2/3}.
$$
Without loss of generality  we may assume that the semi-major axis $a=1$ and 
then set
$x=1-b^2$, so that the entire range of eccentricity is represented by $0< 
x<1$.  The explicit
statement of (Q1) as  posed by M. Vuorinen  is given by the problem: 

\vspace{6pt}
{\em Is it true for  $0<x<1$ that we have}
$$
 F\left(\frac{1}{2}, -\frac{1}{2};1:x \right) - 
\left(\frac{1+(1-x)^{3/4}}{2}
\right)^{2/3} >0?
$$

Let us  denote the error between the normalized elliptical perimeter and the 
algebraic approximation  by
\begin{equation} {\cal E}(x)= F\left(\frac{1}{2},-\frac{1}{2};1:x\right) 
- A(x)
\end{equation} 
where for Muir's approximation, 
$$
 A(x)=\left(\frac{1+(1-x)^{\frac{3}{4}}}{2}  \right)^{2/3}.
$$ 
In the course of answering  (Q1) we will in fact show the stronger result 
that
\begin{equation} .00006 < \frac{{\cal E}(x)}{x^4} < {\cal E}(1) <  
.00666, \hspace{.25in}  0< x
< 1.
\end{equation}
This bound on the error illustrates the surprising  accuracy of Muir's 
approximation.

It is natural to ask whether an estimate of this type  can be obtained 
for each
of the approximations in Table 1. The positivity of the first term in the 
Taylor expansion  of
$L(a,b)/ (\pi (a+b) ) - A(\lambda)$ is consistent with the conjecture 
that Muir's approximation
to $L(a,b)$ could indeed be an approximation from below. Furthermore, the 
power of the first
nonzero term in the Taylor expansion suggests that one consider ${\cal 
E}(x)/ x^4$ in proving the 
result in (2).   For those approximations in Table 1 for which the first 
term in the Taylor expansion
is positive, i.e., for the  approximations of Kepler, Lindner, 
Ramanujan(a,b), Bronshtein, Selmer(a,b)
and Jacobsen, it is natural to ask whether they  are approximations from 
below and if a result
analogous to (2) can be obtained. What we will show in this article  is 
that this is indeed the case. 

An obvious analogy would suggest that the results due to Euler, Sipos, 
Peano,  and Almkvist
are approximations from above.  Somewhat surprisingly, the techniques 
developed in this article are not
directly applicable to these approximations.  However, using  deeper 
methods, it has been
shown  by Barnard, Pearce,and Richards, \cite{KD}, that even stronger 
results can be obtained which
imply that  inequalities analogous to (2) hold, and hence the 
approximations of Euler, Sipos, Peano, 
and Almkvist are indeed from above.

The verification of (2) is obtained by showing the positivity of an 
infinite series in the form 
$$ 
{\cal E}(x)=\sum_{n=0}^{\infty}  b_n(x). 
$$
We first verify that there exists an $N$ for which $b_n(x)>0, n\geq N$. 
Consequently, the proof of
(2), and hence the resolution of (Q1), is reduced  to examining the 
positivity of the polynomial 
$$  
p(x)=\sum_{n=0}^{N}  b_n(x). 
$$  
Because the coefficients of the polynomial $p$ are integers, the 
positivity of $p$ may be verified
by a Sturm sequence argument that involves only exact arithmetic.  In a 
subsequent section, we will
elaborate on the details of this approach but it should be emphasized at 
this point that  these 
calculations, though computationally formidable, are readily carried out 
using a computer algebra
system.  Nevertheless, care must be exercised to ensure that only  exact 
arithmetic and algebra are
performed. 

\vspace{18pt} 
\noindent{\bf{3. Muir's Approximation and Verification of 
Vuorinen's Conjecture}}
\vspace{12pt} 

In this section we will verify (2) and hence answer (Q1). The arguments 
needed for Muir's
approximation turn out to be exceptional among the list of approximations 
in Table 1.  Nevertheless,
the arguments used in this case suggest a more general approach to the  
verification of inequalities
of the type (2) that are relevant to the other approximations.

In order to avoid the complications associated with a choice of a branch 
cut, we first make the
substitution $x \rightarrow 1-x^4$  and define
\begin{equation} 
{\cal G}(x)= \displaystyle{\frac{{\cal E}(1-x^4)}{(1-x^4) ^4}}
\end{equation}
The result is established by showing that $G'(x)$ does not change sign 
and so $G$ is monotone. More
precisely, we will show
$G'(x)<0, \; 0<x<1$.  It is convenient to calculate
$G'$ using a computer algebra system (CAS) such as Mathematica or Maple.  
Using Maple,  we obtain 
$$ 
{\cal G}'(x)= \frac{N1+ N2+P}{D} 
$$ 
where
$$
 N1 = -4\,x \; \sqrt[3]{\frac{1}{2}+\frac{1}{2}{x}^{3}} \; {\it 
K}(\sqrt{1-{x}^{4}})
$$
$$
N2= -28\, x\; \sqrt[3]{\frac{1}{2}+\frac{1}{2}{x}^{3}}\; {\it 
E}(\sqrt{1-{x}^{4}})
$$
$$ 
P= \pi(1+ 7{x}^{4}+8x)
$$
$$
 D =\frac{ \pi \,\left 
(-1+{x}^{4}\right)^{5}\sqrt[3]{\frac{1}{2}+\frac{1}{2}{x}^{3}}}{x^2}
$$ 
and ${\it K}, {\it E}$ denote elliptical integrals of the first and 
second
kind respectively.  It should be noted that the the
above expression for ${\cal G}'(x)$ can be  checked by recalling the  
definitions
of $ F\left(\frac{1}{2},-\frac{1}{2};1:x\right), {\it K}$ and  ${\it 
E}$.

Since  $D<0$, to verify that $ \mathcal{ G}'(x) < 0 $ it suffices to show 
$$ 
 {\cal H}(x)= \frac{-N1-N2}{P} \leq 1.
$$  
 Since ${\cal H}(1)=1$, we need only verify that
${\cal H}'(x) \geq 0 $.  After converting the elliptical integrals  to  
hypergeometric
functions and simplifying,  the problem is reduced to showing
\begin{equation}
p_1(x)\, F\left(\frac{1}{2},\frac{1}{2};1:1-x^4\right) + p_2(x)\, 
F\left(\frac{1}{2},
-\frac{1}{2};1:1-x^4\right)  \geq 0
\end{equation}
where
\begin{align*}
\begin{split}  p_1(x) & =
252{x}^{9} -504{x}^{8} + 756{x}^{7} -502{x}^{6}+ 
248{x}^{5}+6{x}^{4}-4{x}^{3}+6{x}^{2} - 4x+2,\\[1ex]
 p_2(x) &= -238{x}^{6} +196{x}^{5}-154{x}^{4}-112{x}^{3}+94{x}^{2} 
-52x+10.
\end{split}
\end{align*} 
A crucial manipulation involves the  conversion of the above expression 
to  series form. One may then
express the left hand side of (4) as
$$ 
p_1(x) \sum_{k=0}^{\infty} \left(\frac{(  \frac{1}{2}) _k }{ k!}\right)^2
(1-x^4)^k  
 \; + \;  p_2(x)  \sum_{k=0}^{\infty} \frac{ ( -\frac{1}{2}) _k ( 
\frac{1}{2})_k}{k!
k!} (1-x^4)^k.
$$ 
These series may be combined as
\begin{equation}
p_1(x) + p_2(x) + \sum_{k=1}^{\infty} \left(\frac{(  \frac{1}{2}) _k }{
k!}\right)^2 \left [ p_1(x)-p_2(x)/(2k-1) \right ] (1-x^4)^k.  
\end{equation}

It is at this point that a Sturm sequence argument is applied in order to 
verify the needed
positivity. More precisely, we shall use this terminology to refer to a 
generalized form of Descartes'
rule of signs which gives the precise number of zeros of a polynomial in 
an interval.  This is
achieved by  utilizing the Euclidean algorithm and a notion of {\em Sturm 
sequences} (\cite{PH},
\cite{G}).   In particular, given a polynomial $P(x)$, and  a real number 
$a$, define the function
$V(a)$ as  the number of sign variations in the numbers $\{ p_0(a), 
p_1(a), \cdots, p_n(a)\}$.  The
{\em Sturm} polynomials $\{ p_0, p_1, \cdots, p_n\}$ are defined by 
$$
\left\{
\begin{array}{l} p_0(x)=P(x)\\ p_1(x)=P'(x)
\end{array}
\right.
$$ 
and for each $ k\geq 2$, $p_k(x)$ is the unique polynomial of degree less 
than that of $ p_{k-1}(x)$
such that 
$$ 
p_{k-2}(x) =q_k(x)p_{k-1}(x)-p_k(x)\
$$  
where $q_k(x)$ is a polynomial.  Sturm's Theorem states that if $P(a) 
\neq 0$ and $P(b) \neq 0$,  then
the number of distinct roots of $P(x)=0$ in the interval $[a,b]$ is 
exactly $V(a)-V(b)$.  In applying
Sturm's theorem to an arbitrary polynomial,  rounding error may affect 
the sign of $p_i(a)$ and
$p_i(b)$  and  hence limit  the usefulness of the procedure.  However, 
for our purposes, the
polynomials to which the method is always applied have integer 
coefficients and with the exact
arithmetic provided by a computer algebra system, accuracy is never 
compromised.  The method is
implemented by calling the procedure {\em sturm} within Maple. This 
procedure  uses Sturm's theorem to
return the number of real roots in the interval (a,b] of a  polynomial 
$P(x)$.

For the series given in (5), a Sturm sequence argument applied to 
$p_1(x)-p_2(x)/(2k-1) $ shows 
that when $k=4$ this polynomial has no zeros.  It then follows that
\begin{equation}
\displaystyle{ p_1(x) -\frac{p_2(x)}{2k-1}   > 0, \;\;\;  k >4} 
\end{equation} 
The positivity of the infinite series will be confirmed if the  
$61^{st}$ degree polynomial with rational coefficients, 
\begin{equation}
\displaystyle{ p_1(x) +p_2(x) +  \sum_{k=1}^{13} \left(\frac{( 
\frac{1}{2}) _k }{ k!}\right)^2 \left [
p_1(x)-\frac{p_2(x)}{2k-1} \right ] (1-x^4)^k  }
\end{equation} 
is positive.  This is readily verified by  another Sturm
sequence argument.

Because of the topic in Vuorinen's article \cite{MV}, the question  he 
posed regarding Muir's
approximation was  expressed in terms of the eccentricity and hence 
inequality (2) was given as a
function of $e$. If the notation of \cite{AB} is adopted, it is  more 
natural  to express the error in
terms of $\lambda= (a-b)/(a+b)$.   What we shall show in the next section 
is that the methods that
have been utilized in verifying (2) can be generalized to obtain 
analogous results. As argued above,
the proof will ultimately reduce to verifying the positivity of an 
infinite series as in (5), and
subsequently,  the positivity of a polynomial analogous to (7).  


\vspace{18pt} 
\noindent{\bf{4. Approximations of Elliptical Perimeter From Below}}
\vspace{12pt} 


With the exception of Muir's,   the algebraic
approximations from below given in Table 1 may be denoted as 
$A(\lambda^2)$.  There is  a polynomial
$\phi(x)$ so that with the change of variable $\lambda^2=\phi(x)$, we 
obtain a rational function 
$f(x)=A( \phi(x))$. In particular,  for  the approximation  of Kepler,
$\phi(x)=1-x^2$, for those of Lindner, Selmer(a), Bronshstein, and 
Jacobsen, $\phi(x)=x$,  for
Ramanujan(a,b), 
$\phi(x)=(4-x^2)$  and  $\phi(x)=(4-x^2)/3$ respectively,  while for 
Selmer(b), $\phi(x)=2(1-x^2)$. Let
$N$ denote the power of the first nonzero term in the Taylor series for 
the difference
$L(a,b)/( \pi (a+b) ) - A(\lambda^2)$.

The result,  analogous to (3), that will be verified  is   
\begin{equation}
\frac{d}{dx} \left(\frac{F(-\frac{1}{2}, -\frac{1}{2}; 
1:\phi(x))-f(x)}{[\phi(x)]^M}\right) \neq 0.
\end{equation}
Here  $M=N/2$,
and the interval on which   (8)  will hold,  is determined by  the change 
of variable 
$\lambda^2=\phi(x)$. 

 
We begin by computing
\begin{align*}
\begin{split}
& 
\frac{d}{dx} (\frac{F(-\frac{1}{2}, -\frac{1}{2}; 
1:\phi(x))-f(x)}{[\phi(x)]^M}) = 
\frac{1}{[\phi(x)]^{M+1}} [Mf(x)\phi'(x)-\phi(x)f'(x)] +  \\[4ex]
&  
\frac{1}{[\phi(x)]^{M+1}}\left[\frac{\phi(x) \phi'(x)}{4} 
\sum_{k=0}^{\infty}
\frac{(( \frac{1}{2})_k)^2 }{(2)_k} \frac{\phi^k(x)}{k!} \right.  \left.-
M\phi'(x)\sum_{k=0}^{\infty}
\left(\frac{(  -\frac{1}{2}) _k }{k!}\right)^2 \phi^k(x)\right]  \\[4ex]
& = D1  + D2. 
\end{split}
\end{align*} 
For  $D2$ we have
\begin{align*}
\begin{split}
&  D2 =\frac{1}{[\phi(x)]^{M+1}}\left[\frac{\phi(x)\phi'(x) 
}{4}-M\phi'(x)\right] + \\[3ex]
& \hspace{.5in} \frac{1}{[\phi(x)]^{M+1}} \sum_{k=1}^{\infty} 
\left(\frac{(
\frac{1}{2})_k}{k!}\right)^2 
\left[\frac{\phi(x)\phi'(x)}{4(k+1)}-\frac{M\phi'(x)}{(2k-1)^2} \right]
\phi(x)^k\\[4ex] 
& \hspace{.25in} =D21+D22
\end{split}
\end{align*} 
A further computation shows that
\begin{align*}
\begin{split}
&  D22 =   \\[3ex] 
& \frac{\phi'(x)}{[\phi(x)]^{M+1}}\left[\frac{\phi(x)  }{4} \left(-1 + 
  \sum_{k=0}^{\infty} \left(\frac{(\frac{1}{2})_k}{k!}\right)^2 
\frac{\phi^k(x) }{(k+1)} \right) \right. 
\left.- \sum_{k=1}^{\infty} \left(\frac{(\frac{1}{2})_k}{k!}\right)^2 
\frac{M\phi^k(x)}{(2k-1)^2}  \right]
\end{split}
\end{align*} 
\begin{align*}
\begin{split}
&\hspace{-6in}\\  
= \frac{\phi'(x)}{[\phi(x)]^{M+1}} \left[ \frac{\phi(x) }{4} \left( -1+ 
   \sum_{k=1}^{\infty} \left( \frac{(\frac{1}{2})_{k-1}}{(k-1)!}\right)^2 
\frac{\phi^{(k-1)}(x) }{k} \right) \right. 
\left. - \sum_{k=1}^{\infty} \left(\frac{(\frac{1}{2})_k}{k!}\right)^2 
\frac{M \phi^k(x)}{(2k-1)^2}  \right]&
\end{split}
\end{align*}
\begin{align*}
\begin{split}
&\hspace{-6in}\\  
= \frac{\phi'(x)}{[\phi(x)]^{M+1}}\left [\frac{-\phi(x) }{4}+ 
   \sum_{k=1}^{\infty} \left(\frac{(\frac{1}{2})_{k}}{k!}\right)^2 
\frac{k[\phi(x)]^{k} }{(2k-1)^2}  \right. 
\left.- \sum_{k=1}^{\infty} \left(\frac{(\frac{1}{2})_k}{k!}\right)^2 
\frac{M\phi^k(x)}{(2k-1)^2}  \right]&
\end{split}
\end{align*} \\

 \noindent  Combining these results gives us 
\begin{align*}
\begin{split} 
&  \frac{d}{dx} (\frac{F(-\frac{1}{2}, -\frac{1}{2}; 
1:\phi(x))-f(x)}{[\phi(x)]^M}) = 
\\[4ex] 
 &  \frac{1}{[\phi(x)]^{M+1}} \left[ M\phi'(x)(f(x)-1)- f'(x)\phi(x) 
  + 
\phi'(x) \sum_{k=1}^{\infty} \left(\frac{(\frac{1}{2})_k}{k!}\right)^2 
\frac{(k-M)\phi^k(x)}{(2k-1)^2}  \right]
\end{split}
\end{align*} 

The  necessary inequality that is analogous to (6) is trivially met if 
$k>M$. To verify (8), one 
need only determine $K \geq M$ and apply a Sturm sequence argument to 
show that the polynomial 
$$
Mf(x)\phi'(x)-f'(x)\phi(x) -M\phi'(x) + \phi'(x) \sum_{k=1}^{K}
\left(\frac{(\frac{1}{2})_k}{k!}\right)^2 
\frac{(k-M)\phi^k(x)}{(2k-1)^2} 
$$
has the same sign as 
$$
\phi'(x) \sum_{k=1}^{\infty} \left(\frac{(\frac{1}{2})_k}{k!}\right)^2 
\frac{(k-M)\phi^k(x)}{(2k-1)^2}.  
$$  

From the monotonicity of  $  \left[ 
F\left(-\frac{1}{2},-\frac{1}{2};1:\phi(x)\right)
-  A(\phi(x))\right]/[\phi(x)]^M$, one readily obtains, analogous to 
inequality (2),  bounds
$(\alpha, \beta)$ on 
$$ 
\frac{{\cal E}(\lambda)}{\lambda^N}= 
\frac{F\left(-\frac{1}{2},-\frac{1}{2};1:\lambda^2\right)
-  A(\lambda^2)}{\lambda^N}. 
$$ 
In particular, $\alpha$ is obtained as the limit of ${\cal 
E}(\lambda)/\lambda^N$
as $\lambda \rightarrow 0$, while
$\beta$ is determined by evaluation at $\lambda=1$.
The following table illustrates the accuracy of the  approximations from 
below by specify $\alpha,
\beta$ where
\begin{equation}
\alpha < \frac{{\cal E}(\lambda)}{\lambda^N} < {\cal E}(1) <  \beta, 
\;\;\; 0<\lambda<1.
\end{equation}

\vspace {12pt} 
\noindent Table 2.
$$
\begin{array}{lc}
& \hspace{.75in} (\alpha,\beta) \\ [1ex]
\hline \hline
\mbox{Kepler (1609)} & \hspace{.96in} \mbox{ (.75,\;1.2732) } \\[1ex]
\mbox{Lindner (1904)} & \hspace{.75in}\mbox{ (.003906,\;.007614) } \\[1ex]
\mbox{Ramanujan (1914 a)} & \hspace{.75in}\mbox{ (.001953,\;.005290) } 
\\[1ex]
\mbox{Ramanujan (1914 b)} & \hspace{.75in}\mbox{ (.000022,\;.000512) } 
\\[1ex]
\mbox{Bronshtein (1964)} & \hspace{.75in}\mbox{ (.002929,\;.006572) 
}\\[1ex]
\mbox{Selmer (1975 a)} & \hspace{.75in}\mbox{ (.000549,\;.002406) } 
\\[1ex]
\mbox{Selmer (1975 b)} &\hspace{.75in} \mbox{ (.000305,\;.001792) } 
\\[1ex]
\mbox{Jacobsen (1978)} & \hspace{.75in}\mbox{ (.000125,\;.001130) }\\[1ex]
\end{array}
$$

 


To illustrate the use of these bounds in obtaining error approximations 
first note that 
$$
\frac{a-b}{a+b}=\frac{1-\sqrt{1-e^2}}{1+\sqrt{1-e^2}}=\lambda(e).
$$
It follows that the error in elliptical perimeter approximation 
$$
\tilde{E}(\lambda):= L(a,b) - \pi (a+b) F(\frac{-1}{2}, \frac{-1}{2}; 1:\lambda^2)
$$
satisfies
\begin{equation} 
a \pi \alpha   (1 + \sqrt{1-e^2}) [\lambda(e)]^N < \tilde{E}(\lambda)<a \pi \beta 
  (1 + \sqrt{1-e^2})
[\lambda(e)]^N
\end{equation} 



                            

\vspace{18pt} 
\noindent{\bf{5. Implications and Further Directions }}
\vspace{12pt} 


It can be shown that the remaining approximations
in Table 1 are from above.   To verify this requires considerably more 
sophisticated methods. In fact,
what has been shown by  Barnard, Pearce and Richards \cite{KD} is that, 
with the exception of
Euler's,  the  coefficients of the Taylor expansions of
${\cal E}(x)$ are all of the same sign.  As a consequence of their 
results one can deduce not only that
(8) holds, but more generally,  for every integer $n \geq 0$, 
$$
\frac{d^n{\cal E}(x)}{dx^n} \neq0.
$$ 

Of course, the results of this
paper are trivial consequences of this more general result,  but the 
methods that are presented
here  are of interest in their own right.  For instance,  the results presented here, and in particular
equation (10),  provides a relatively simple derivation of estimates that historically appeared in the context
of  astronomical considerations  and the application of Landen's transformation
 \cite{AB}. 




In addition, there is an obvious intriguing issue associated with  the study of these 
inequalities.
Is there a geometrical significance to these algebraic approximations that
accounts for the monotone properties that these approximations possess?  
It has been suggested by
Berndt \cite{B} that Ramanujan's formulas have their origins  in the 
representation of analytic
functions as continued fractions.  As for the other approximations, 
little is known as to
their motivation.  


\begin{thebibliography}{99}

\bibitem{AB}  Almkvist, G., Berndt, B., ``Gauss,Landen, Ramanujan, the
arithmetic-geometric mean, ellipses, $\pi$, and the {\em Ladies Diary},'' 
 {\em The
American Mathematical Monthly}, 95,1988 (585-608).
\bibitem{B} Berndt, B., {\em Ramanujan's Notebooks, Part III}, Springer Verlag 
(1985). 
\bibitem{KD}  Barnard, R.W., Pearce, K. and Richards, K., ``An inequality 
involving
the generalized hypergeometric function and the arc length of an 
ellipse,'' 
preprint. 
\bibitem{KD}  Devlin, K., ``The logical structure of computer-aided 
mathematical
reasoning,''  {\em The American Mathematical Monthly},104,1997 (632-646). 
\bibitem{PH}  Henrici, P., {\em Applied and Computational Complex Analysis}, 
John Wiley \&
Sons,  New York (1974)
\bibitem{MV}  Vuorinen, M., ``Hypergeometric functions in geometric  
function
theory,''  {\em Proceedings of Special Functions and Differential 
Equations}, Eds.: K.
Srinivasa Roa, R Jagannthan, G. Van der Jeugy, Allied Publishers (1998)
\bibitem{G}  Young, D. and Gregory, R., {\em A Survey of Numerical 
Mathematics}, Dover, 
New York (1973)


\end{thebibliography}


\end{document}

